\section{Experimental Methodology}

We investigate the proposed vision-based rendezvous architecture with a set of controlled numerical experiments. Each main element is independently tested before investigating the full navigation-and-control pipeline. We move from analytically validating the dynamics through to estimating the visual poses, robust estimation, recursive state estimation, guidance-and-control, and lastly end-to-end Monte Carlo testing. 

This allows us to examine any numerical and algorithmic errors separate from the effects at the system level. 

Specifically, we demonstrate validity of the dynamics model against a known analytic solution and then characterising the pose-estimation pipeline’s sensitivity to image noise and outliers. We then examine the state estimators using a controlled measurement model and the fullvision--estimation--control pipeline using simulate camera measurements. All experiments use the same parameters of the target-spacecraft and the camera model introduced in sectionIII unless stated otherwise. A random seed will be chosen that will remain constant within a single Monte Carlo trial to allow for reproduceability.

\subsection{Experimental Objectives}

The experiment campaign is structured around six main investigations.

\begin{enumerate}

\item \textbf{Dynamics validation:} The numerical implementation of relative dynamics is verified against the analytical HCW state transition solution, and numerical solver convergence is analyzed.

\item \textbf{Pixel-noise sensitivity:} The image plane measurement noise level is varied and its impact on EPnP pose is analyzed along with the effectiveness of the nonlinear pose refinement stage.

\item \textbf{Outlier robustness:} RANSAC relative pose is analyzed under an increasing amount of outliers.

\item \textbf{State-estimation comparison:} Under an assumed Gaussian measurement model, the MEKF and UKF stat estimation algorithms are compared.

\item \textbf{Guidance and control comparison:} Relative rendezvous under HCW dynamics is assessed for guidance and control, using a receding-horizon MPC as well as a LQR controller.

\item \textbf{End-to-end robustness:} The overall vision to control architecture performance is assessed for various values of pixel noise based on Monte Carlo simulations.

\end{enumerate}

Consequently the experiments can be structured to address three escalating and broader issues : that the constituent mathematical models hold , that individual navigation and control algorithms perform in presence of unwanted perturbations , and finally that the overall architecture continues to function correctly when individual components interact with each other.

\subsection{Simulation Environment and Common Parameters}

All simulation work is conducted on a numerical simulation infrastructure developed in Python. Relative translational dynamics are modeled in the LVLH reference frame with the HCW equations while target motion in attitude is propagated by a built-in rigid body attitude dynamics model. Numerical propagation is performed with standard fourth-order Runge-Kutta integration in modules that require it.

A parametric Ariane 44L upper stage geometry serves as the target spacecraft, containing 15 analytical keypoints in its body frame coordinates; which are then projected to the imaging sensors with the use of a calibrated camera model given in Section III-D.

Experiments with full images use $1024 \times 1024$ pixel frames with noise added to the ideal projected keypoints through a zero-mean Gaussian. Unless otherwise stated standard pixel noise level of $\sigma_{\mathrm{px}}=1.5$ px is utilized for all standard vision experiments.

Table~\ref{tab:experimental_parameters} summarises the standard simulation parameters used throughout the experimental campaign.
\begin{table}[t]
\centering
\caption{Common simulation and experimental parameters.}
\label{tab:experimental_parameters}
\begin{tabular}{lc}
\toprule
\textbf{Parameter} & \textbf{Value} \\
\midrule
Camera image size & $1024\times1024$ px \\
Target keypoints & 15 \\
Nominal pixel noise & $1.5$ px \\
HCW initial range (Exp. D) & 50 m \\
Filter sampling interval & 10 s \\
Filter horizon (Exp. D) & 3000 s \\
MEKF/UKF position noise & 0.5 m \\
MEKF/UKF attitude noise & 0.05 rad \\
LQR/MPC control interval & 10 s \\
MPC prediction horizon & 20 steps \\
Maximum control input (Exp. E) & 0.1 m/s$^2$ \\
Maximum control input (Exp. F) & 0.3 m/s$^2$ \\
\bottomrule
\end{tabular}
\end{table}

\subsection{Dynamics Validation}

The relative translational dynamics are first validated independently of
the vision and estimation subsystems. Experiment A compares the
closed-form HCW state-transition formulation against numerical
integration using the DOP853 solver.

A periodic HCW trajectory with an initial radial amplitude of 100 m is
propagated for three orbital periods. The analytical state is evaluated
at the same time points as the numerical solution, and the position and
velocity discrepancies are computed as

\begin{equation}
    e_r(t)
    =
    \left\|
        \mathbf{r}_{\mathrm{analytical}}(t)
        -
        \mathbf{r}_{\mathrm{numerical}}(t)
    \right\|_2,
\end{equation}

and

\begin{equation}
    e_v(t)
    =
    \left\|
        \mathbf{v}_{\mathrm{analytical}}(t)
        -
        \mathbf{v}_{\mathrm{numerical}}(t)
    \right\|_2.
\end{equation}

The maximum errors over the complete propagation interval are used as
the primary validation metrics.

The numerical integration is initially performed using
$\mathrm{rtol}=10^{-12}$ and $\mathrm{atol}=10^{-14}$. A second
convergence experiment varies the relative tolerance over

\begin{equation}
    \mathrm{rtol}
    \in
    \left\{
    10^{-6},10^{-8},10^{-10},10^{-12},10^{-13}
    \right\},
\end{equation}

with $\mathrm{atol}=10^{-2}\mathrm{rtol}$. The resulting maximum
position and velocity discrepancies are used to verify convergence of
the numerical propagation.

\subsection{Pixel-Noise Sensitivity of Pose Estimation}

Experiment B evaluates the sensitivity of the visual pose-estimation
pipeline to image-plane measurement noise. The ideal two-dimensional
keypoint observations are corrupted using independent zero-mean
Gaussian noise,

\begin{equation}
    \tilde{\mathbf{u}}_i
    =
    \mathbf{u}_i
    +
    \boldsymbol{\eta}_i,
    \qquad
    \boldsymbol{\eta}_i
    \sim
    \mathcal{N}
    \left(
        \mathbf{0},
        \sigma_{\mathrm{px}}^2\mathbf{I}_2
    \right).
\end{equation}

The experiment considers

\begin{equation}
\sigma_{\mathrm{px}}
\in
\{0.1,\;0.25,\;0.5,\;1,\;1.5,\;2,\;3,\;5\}\ {\rm px}.
\end{equation}

For each noise level, 200 independent Monte Carlo trials are performed.
The raw EPnP solution is compared with the solution obtained after
nonlinear Gauss--Newton pose refinement.

The translation error is evaluated using the Euclidean norm,

\begin{equation}
    e_t =
    \left\|
        \hat{\mathbf{t}}-\mathbf{t}_{\mathrm{gt}}
    \right\|_2,
\end{equation}

while the rotational error is evaluated using the geodesic angle

\begin{equation}
    e_R
    =
    \cos^{-1}
    \left[
    \frac{
    \operatorname{tr}
    \left(
    \hat{\mathbf{R}}
    \mathbf{R}_{\mathrm{gt}}^{T}
    \right)-1
    }{2}
    \right].
\end{equation}

The mean reprojection error is additionally recorded to quantify the
image-space consistency of each recovered pose. Translation RMSE,
rotation error, and reprojection error are reported as functions of
pixel noise.

\subsection{Robustness to Image Outliers}

Experiment C evaluates the robustness of pose estimation to incorrect
keypoint correspondences. Inlier observations are corrupted with
Gaussian pixel noise of $\sigma_{\mathrm{px}}=1$ px. A controlled
fraction of the keypoints is then replaced with randomly generated
image coordinates distributed over the complete camera image.

The tested outlier fractions are

\begin{equation}
    f_{\mathrm{out}}
    \in
    \{0,\;5,\;10,\;20,\;30\}\%.
\end{equation}

For each condition, 200 independent trials are performed. Bare EPnP is
compared with the proposed RANSAC--EPnP configuration using a
2-pixel inlier threshold and a maximum of 500 RANSAC iterations.

A trial is classified as a failure if either the pose solver reports an
invalid solution or the recovered rotation error exceeds $15^\circ$.
The resulting failure rate is

\begin{equation}
    P_{\mathrm{fail}}
    =
    \frac{N_{\mathrm{invalid}}
    +
    N_{e_R>15^\circ}}
    {N_{\mathrm{trials}}}
    \times 100\%.
\end{equation}

Translation and rotation RMSE are computed over successful trials,
while the failure rate is computed over all trials. This distinction
prevents unsuccessful trials from being silently removed from the
robustness analysis.

\subsection{State-Estimation Evaluation}

Experiment D evaluates the recursive state-estimation stage by
comparing the MEKF and UKF under an identical controlled measurement
model. The purpose of this experiment is to isolate filter behavior
from the stochastic characteristics of the camera and PnP pipeline.

The simulated rendezvous begins at a 50-m along-track separation with
zero initial relative velocity. The target is assigned a slow
uncontrolled angular velocity of

\begin{equation}
    \boldsymbol{\omega}
    =
    [0.01,\;0.02,\;-0.015]^{T}\ {\rm rad/s}.
\end{equation}

The filter sampling interval is 10 s and the simulation contains 300
steps, corresponding to 3000 s.

The initial filter state contains deliberate errors of approximately
4 m in position and $17^\circ$ in attitude to represent a cold-start
condition. Both filters receive the same noisy pose measurement with

\begin{equation}
    \sigma_{\mathrm{pos}}=0.5~{\rm m},
    \qquad
    \sigma_{\mathrm{att}}=0.05~{\rm rad}.
\end{equation}

Position, attitude, and velocity errors are recorded throughout the
simulation. The root-mean-square error is computed as

\begin{equation}
    \mathrm{RMSE}
    =
    \sqrt{
    \frac{1}{N}
    \sum_{k=1}^{N}
    e_k^2
    }.
\end{equation}

Maximum position and attitude errors, the number of accepted
measurements, and wall-clock execution time are also recorded.

Importantly, this experiment uses direct synthetic Gaussian pose
measurements rather than pixel measurements generated by the camera
and PnP stages. Thus, the results characterize the filter comparison
under a controlled measurement distribution and should not be
interpreted as an end-to-end camera-navigation comparison.

\subsection{Guidance and Control Evaluation}

Experiment E compares an infinite-horizon discrete LQR controller with
a receding-horizon MPC controller for translational rendezvous under
HCW dynamics.

Both controllers start from the same relative state,

\begin{equation}
    \mathbf{x}_0
    =
    [0,\;50,\;0,\;0,\;0,\;0]^T,
\end{equation}

corresponding to a 50-m along-track separation and zero initial
relative velocity. The control interval is 10 s and the simulation
contains 200 control steps.

The LQR controller is obtained from the discrete algebraic Riccati
equation \cite{anderson1990}, while MPC \cite{mayne2000} uses a
20-step prediction horizon and the LQR solution as its terminal cost,
following the approach of Breger and How~\cite{breger2008} for
spacecraft proximity operations. Both controllers use identical state
and control weighting and a maximum control input magnitude of
$0.1~\mathrm{m/s^2}$.

The controllers are compared using final range,

\begin{equation}
    r_f = \|\mathbf{r}_f\|_2,
\end{equation}

total velocity increment,

\begin{equation}
    \Delta v
    =
    \sum_{k=0}^{N-1}
    \|\mathbf{u}_k\|_2\Delta t,
\end{equation}

maximum control magnitude, and computational cost per control step.

For this controlled comparison, both controllers are evaluated without
the optional approach-cone constraint so that differences are not
caused by constraint feasibility.

\subsection{End-to-End Monte Carlo Evaluation}

Experiment F evaluates the complete rendezvous architecture by
combining relative dynamics, visual measurements, state estimation,
and guidance and control. Two complementary studies are performed.

\subsubsection{Monte Carlo Configuration Study}

The first study performs 25 Monte Carlo trials for each of three
configurations:

\begin{enumerate}
    \item \textbf{Perfect-state configuration:} ideal state knowledge
    with no measurement noise and no initial-state error.

    \item \textbf{EKF-only configuration:} direct noisy state
    measurements are supplied to the EKF, without the camera and
    visual pose-estimation stages.

    \item \textbf{Full-pipeline configuration:} simulated camera
    measurements are processed through the visual pose-estimation
    pipeline and subsequently supplied to the EKF and LQR controller.
\end{enumerate}

Each trial contains 250 simulation steps. The initial relative position
is

\begin{equation}
    \mathbf{r}_0
    =
    [30,\;3,\;-1.5]^T~{\rm m},
\end{equation}

and the initial relative velocity is

\begin{equation}
    \mathbf{v}_0
    =
    [-0.08,\;0,\;0]^T~{\rm m/s}.
\end{equation}

For the noisy configurations, the initial position error is

\begin{equation}
    \delta\mathbf{r}_0
    =
    [2.5,\;-0.8,\;0.5]^T~{\rm m}.
\end{equation}

The nominal full-vision configuration uses
$\sigma_{\mathrm{px}}=1.5$ px.

The primary metrics are final range, total $\Delta v$, and docking
success rate. These metrics quantify both navigation/control
performance and mission-level outcome.

\subsubsection{Pixel-Noise Ablation}

A second Monte Carlo study evaluates the sensitivity of the complete
vision-based pipeline to image measurement noise. The full
vision--EKF--LQR configuration is evaluated at

\begin{equation}
    \sigma_{\mathrm{px}}
    \in
    \{0.3,\;0.7,\;1.0,\;1.5,\;2.0,\;3.0,\;4.0\}\ {\rm px}.
\end{equation}

Twenty independent trials are performed at each noise level. The mean
and standard deviation of the final range are recorded as

\begin{equation}
    \bar{r}_f
    =
    \frac{1}{N}
    \sum_{i=1}^{N}r_{f,i},
\end{equation}

\begin{equation}
    \sigma_{r_f}
    =
    \sqrt{
    \frac{1}{N}
    \sum_{i=1}^{N}
    (r_{f,i}-\bar{r}_f)^2
    }.
\end{equation}

Docking success is additionally reported as a function of pixel noise.
This experiment provides a system-level sensitivity analysis linking
the quality of the visual measurements to the final rendezvous
performance.

\subsection{Evaluation Metrics}

The experimental campaign uses metrics appropriate to each stage of the
rendezvous architecture. Dynamics validation is evaluated using maximum
position and velocity propagation errors. Visual pose estimation is
evaluated using translation error, rotation error, and image-plane
reprojection error. Robustness is quantified using solver failure rate
and estimation RMSE.

For state estimation, position and velocity errors are computed using
the Euclidean norm,

\begin{equation}
    e_{\mathbf{r},k}
    =
    \|\hat{\mathbf{r}}_k-\mathbf{r}_k\|_2,
\end{equation}

\begin{equation}
    e_{\mathbf{v},k}
    =
    \|\hat{\mathbf{v}}_k-\mathbf{v}_k\|_2.
\end{equation}

The attitude error is represented by the geodesic angle between the
estimated and true orientations,

\begin{equation}
    e_{\mathbf{R},k}
    =
    \cos^{-1}
    \left(
    \frac{
    \operatorname{tr}
    (\hat{\mathbf{R}}_k\mathbf{R}_k^T)-1
    }{2}
    \right).
\end{equation}

For a sequence of $N$ samples, the corresponding RMSE is

\begin{equation}
    \mathrm{RMSE}(e)
    =
    \sqrt{
    \frac{1}{N}
    \sum_{k=1}^{N}e_k^2
    }.
\end{equation}

For the control experiments, total $\Delta v$ is used as a measure of
control effort, while final range measures rendezvous accuracy.
Computational cost is reported using wall-clock execution time or
average computation time per control step, depending on the experiment.

For the end-to-end Monte Carlo experiments, docking success is reported
as the fraction of trials satisfying the implemented docking condition
within the simulation horizon.

